% ============================================================
% IV. RESULTADOS
% ============================================================

\section{Resultados}

Los resultados obtenidos mediante el procesamiento de las señales se resumen en la Tabla~\ref{tab:resultados}.

\begin{table}[!t]
\centering
\caption{Resultados del análisis de las señales (ventana de 3 s)}
\label{tab:resultados}
\small
\begin{tabular}{@{}l l r r r@{}}
\toprule
Categoría & Señal & Media & Desv. Est. & $f_d$ (Hz) \\
\midrule

Animal & Gato & 0.036 & 2288.47 & 1940.90 \\
Animal & Perro & -0.107 & 2047.43 & 491.00 \\
Animal & Gallo & 0.049 & 3006.66 & 1334.90 \\

\midrule

Instrumento & Contrabajo & -1.244 & 3836.46 & 146.48 \\
Instrumento & Piano & -0.101 & 1745.00 & 384.69 \\
Instrumento & Flauta & 0.585 & 3826.29 & 432.00 \\

\midrule

Voz & Persona 1 & 6.699 & 2778.54 & 214.30 \\
Voz & Persona 2 & -0.167 & 2153.47 & 142.47 \\

\bottomrule
\end{tabular}
\end{table}

\subsection{Resultados de los sonidos animales}

En los sonidos animales se obtuvieron frecuencias dominantes de $1940.90$ Hz para el gato, $491.00$ Hz para el perro y $1334.90$ Hz para el gallo. La diferencia entre las frecuencias dominantes permitió observar, en términos absolutos, una separación espectral entre las tres muestras: el perro presentó la frecuencia dominante más baja, mientras que el gato presentó la más alta. La comparación de los espectros puede observarse en la Fig.~\ref{fig:animales_comparacion}.

\subsection{Resultados de los instrumentos}

Para los instrumentos se obtuvieron frecuencias dominantes de $146.48$ Hz para el contrabajo, $384.69$ Hz para el piano y $432.00$ Hz para la flauta. El contrabajo presentó claramente la frecuencia dominante más baja; el piano y la flauta presentaron valores superiores, aunque relativamente cercanos entre sí. Es importante considerar que los instrumentos musicales no producen una única frecuencia, sino una frecuencia fundamental acompañada de componentes armónicos, por lo que el espectro completo (Fig.~\ref{fig:instrumentos_comparacion}) proporciona información adicional a la frecuencia dominante aislada.

\subsection{Resultados de las voces}

Las frecuencias dominantes obtenidas para las dos personas fueron $214.30$ Hz y $142.47$ Hz, respectivamente. Aunque ambas personas pronunciaron la misma frase, se observaron diferencias en sus espectros frecuenciales (Fig.~\ref{fig:personas_comparacion}). Una sola grabación por persona no es suficiente para establecer que estas características permitan identificar de manera general a un individuo.

\subsection{Separabilidad estadística cuantitativa}
\label{sec:resultados-separabilidad}

Aplicando la Ec.~\ref{eq:separabilidad} a cada par de señales dentro de una misma categoría se obtiene la Tabla~\ref{tab:separabilidad}, generada con la misma lógica que implementa \texttt{compare\_categories}.

\begin{table}[!t]
\centering
\caption{Puntaje de separabilidad estadística por pares}
\label{tab:separabilidad}
\small
\begin{tabular}{@{}l r r r l@{}}
\toprule
Par & $\Delta f_d$ (Hz) & $\sigma_a+\sigma_b$ & $S_{a,b}$ & Veredicto \\
\midrule

Gato--Perro & 1449.90 & 4335.90 & 0.33 & Solapados \\
Gato--Gallo & 606.00 & 5295.13 & 0.11 & Solapados \\
Perro--Gallo & 843.90 & 5054.09 & 0.17 & Solapados \\

\midrule

Contrabajo--Piano & 238.21 & 5581.46 & 0.04 & Solapados \\
Contrabajo--Flauta & 285.52 & 7662.75 & 0.04 & Solapados \\
Piano--Flauta & 47.31 & 5571.29 & 0.01 & Solapados \\

\midrule

Persona 1--Persona 2 & 71.83 & 4932.01 & 0.01 & Solapados \\

\bottomrule
\end{tabular}
\end{table}

Con el umbral definido en la Sección~\ref{sec:separabilidad} ($S_{a,b} \geq 1$ para separabilidad parcial), \emph{ningún} par resulta siquiera parcialmente separable una vez que la diferencia de frecuencia dominante se normaliza por la dispersión de amplitud combinada: los puntajes van de $0.01$ a $0.33$. Esto contrasta con la lectura puramente cualitativa de la Sección anterior, donde las frecuencias dominantes por sí solas sugerían una separación más clara entre categorías (por ejemplo, Gato vs. Perro).
